% =====================================================================
% Mục 0 — Tóm tắt  (BẢN TIẾNG VIỆT)
%
% VIET SAU CUNG, theo dung yeu cau cua ke hoach: chi soan muc nay khi
% Muc 4, Muc 5, Muc 6 va toan bo bang/hinh da dong bang.
%
% RANG BUOC BAT BUOC: khong duoc dua ra ket qua chua xuat hien trong Muc 4.
% Moi con so o day deu phai co mat trong sections_vi/04_results.tex.
% Kiem tra tu dong bang scripts/person3/check_vi_plan.py.
%
% Cau truc theo mo hinh 5 thanh phan cua skill academic-paper
% (references/abstract_writing_guide.md), do dai 150-250 tu theo
% scientific-writing.md:
%   1. Boi canh va khoang trong   (2 cau)
%   2. Muc tieu                   (1 cau)
%   3. Phuong phap                (2 cau)
%   4. Ket qua chinh              (4 cau, co so lieu cu the)
%   5. Y nghia                    (2 cau)
%
% Khong trich dan tai lieu trong tom tat. Khong dung dau gach ngang lam
% dau cau. Khong dung viet tat chua giai nghia.
%
% Doi chieu so lieu voi Muc 4:
%   99.561 +/- 0.354 ; 92.356 +/- 1.022 ; 98.371 +/- 0.319   Bang 1
%   4.22 M so voi 27.83 M                                    Bang 5
%   61 trong 204 loi                                         Muc 4.3
%   95.3 / 96.6 / 37.970                                     Bang 4
%   17 loai, 532 anh test, 3 seed                            Muc 4.1
% =====================================================================

\begin{abstract}
Ảnh mặt cắt ngang ở thang đo vĩ mô là hướng khả thi nhất để định danh gỗ cây lá
kim bên ngoài phòng thí nghiệm. Chưa rõ bài toán này đòi hỏi kiến trúc học sâu
nào, và bao nhiêu phần trong độ chính xác đã công bố phản ánh bằng chứng giải
phẫu.

Nghiên cứu đánh giá đối sánh năm kiến trúc trên cùng một giao thức và kiểm tra
xem độ chính xác đó dựa trên loại thông tin ảnh nào. ResNet-50,
EfficientNet-B3, MobileNetV3-Large, ConvNeXt-Tiny và Swin-Tiny được tinh chỉnh
trên 17 loài cây lá kim, với phân hoạch gom theo khóa nguồn, ba seed mỗi cấu
hình và một registry đánh giá bất biến. Ba phép thăm dò phá hủy cấu trúc giải
phẫu được huấn luyện trên cùng phân hoạch.

ConvNeXt-Tiny đạt $99.561 \pm 0.354$\,\% độ chính xác test, còn ResNet-50 chỉ
đạt $92.356 \pm 1.022$\,\%. MobileNetV3-Large đạt $98.371 \pm 0.319$\,\% với
4.22\,M tham số so với 27.83\,M của mô hình dẫn đầu. Lỗi mang tính tập trung: 61 trong
tổng số 204 lỗi là cùng một nhầm lẫn giữa \textit{Pinus krempfii} và
\textit{Pinus dalatensis}, cặp duy nhất xuất hiện ở cả năm kiến trúc. Mô hình
huấn luyện trên đầu vào $16\times16$ vẫn giữ 95.3\,\% độ chính xác của ảnh đầy
đủ, ảnh làm mờ nặng giữ 96.6\,\%, còn mô hình chỉ dùng histogram màu đạt
37.970\,\%.

Cấu trúc giải phẫu không sống sót qua những phép biến đổi đó. Các mức độ chính
xác ở đây vì vậy mô tả hiệu năng trên bộ dữ liệu này chứ không phải khả năng
định danh theo giải phẫu, và kiểm chứng trên ảnh thu thập độc lập là bước tiếp
theo bắt buộc.
\end{abstract}

\begin{paperkeywords}
định danh gỗ, loài cây lá kim, ảnh mặt cắt ngang vĩ mô, đánh giá đối sánh học
sâu, học theo đường tắt, rò rỉ dữ liệu
\end{paperkeywords}
